\chapter{Propuesta}


En este capítulo se presenta la propuesta de solución diseñada para automatizar la extracción y validación de información contenida en documentos financieros y administrativos. Se describe el propósito, los tipos de documentos que procesa, su arquitectura general y la forma en que se integra dentro de un flujo operativo. De este modo, se busca mostrar cómo la propuesta responde a la problemática identificada y por qué constituye una alternativa viable, escalable y útil para entornos \textit{fintech}.


\section{Solución propuesta}
La propuesta se desarrolla para automatizar la extracción, estructuración y validación de la información contenida en documentos financieros y administrativos en formatos digitales, tales como PDF e imágenes. Su alcance comprende documentos utilizados en procesos de validación, \textit{compliance} y tesorería, entre ellos estados de cuenta, constancias de situación fiscal, comprobantes de domicilio, documentos de identidad y otros soportes relevantes para la operación. La información extraída constituye el \textit{output} que recibe el usuario del sistema, presentada de manera estandarizada y lista para su utilización en procesos de verificación y control. De este modo, el sistema transforma documentos no estructurados en información útil y operable para la toma de decisiones.

El diseño adopta un enfoque modular que integra distintas arquitecturas de extracción, incluyendo métodos determinísticos basados en expresiones regulares, modelos de lenguaje y enfoques híbridos que combinan ambos principios. Esta diversidad permite comparar estrategias, medir su desempeño y determinar con mayor rigor cuál ofrece el mejor equilibrio entre precisión, robustez y estabilidad. Como resultado del proceso comparativo, la implementación prioriza la arquitectura que presenta el mejor desempeño en función de los criterios definidos de evaluación. Ello permite sustentar técnicamente la selección del método final y asegurar que la alternativa implementada responda a criterios objetivos de calidad y confiabilidad.


% Paleta del proyecto:
%   Entrada/datos:    fill=gray!8,    draw=gray!60
%   Proceso/eval:     fill=orange!8,  draw=orange!50
%   Resultado/output: fill=teal!8,    draw=teal!50
%   General/etapa:    fill=cyan!10,   draw=black
\begin{figure}[H]
\centering
\resizebox{0.75\textwidth}{!}{
\begin{tikzpicture}[
    node distance=0.8cm and 0.6cm,
    every node/.style={font=\small, align=center},
    entrada/.style={
        rectangle,
        rounded corners=6pt,
        draw=gray!60,
        thick,
        minimum width=4cm,
        minimum height=1.8cm,
        fill=gray!8
    },
    proceso/.style={
        rectangle,
        rounded corners=6pt,
        draw=orange!50,
        thick,
        minimum width=7cm,
        minimum height=1.8cm,
        fill=orange!8
    },
    resultado/.style={
        rectangle,
        rounded corners=6pt,
        draw=teal!50,
        thick,
        minimum width=7cm,
        minimum height=1.8cm,
        fill=teal!8
    },
    flecha/.style={
        -{Latex[length=3mm]},
        thick
    }
]

% Métodos (fila superior)
\node[entrada] (det) {
    Métodos\\determinísticos\\[2pt]
    {\scriptsize basados en expresiones regulares}
};
\node[entrada, right=0.8cm of det] (llm) {
    Modelos de\\lenguaje\\[2pt]
    {\scriptsize (\textit{LLMs})}
};
\node[entrada, right=0.8cm of llm] (hyb) {
    Enfoques\\híbridos\\[2pt]
    {\scriptsize \textit{LLM} + determinístico}
};

% Evaluación comparativa
\node[proceso, below=2cm of llm] (eval) {
    Evaluación comparativa\\[2pt]
    {\scriptsize Precisión, robustez y estabilidad}
};

% Arquitectura seleccionada
\node[resultado, below=1cm of eval] (arq) {
    Arquitectura seleccionada\\[2pt]
    {\scriptsize para la versión final del producto}
};

% Flechas convergentes
\draw[flecha] (det.south) |- ([yshift=0.6cm]eval.north) -- (eval.north);
\draw[flecha] (llm.south) -- ([yshift=0.6cm]eval.north) -- (eval.north);
\draw[flecha] (hyb.south) |- ([yshift=0.6cm]eval.north) -- (eval.north);

% Flecha de evaluación a arquitectura
\draw[flecha] (eval.south) -- (arq.north);

\end{tikzpicture}
}
\caption{Esquema del proceso de evaluación y selección de la arquitectura de extracción final.}
\label{fig:seleccion_arquitectura}
\end{figure}

El producto se ofrece a la \textit{startup fintech} cliente como una API REST que puede integrarse directamente en sus propias plataformas y flujos operativos. A través de esta interfaz, el cliente puede consumir el servicio de extracción de documentos bajo demanda, enviando el documento en formato PDF, JPG o PNG al \textit{endpoint} correspondiente para que el sistema procese la solicitud y devuelva los campos extraídos en formato JSON. De esta manera, la solución se incorpora de forma flexible en procesos reales de validación, \textit{onboarding} o \textit{compliance}, sin requerir que el usuario interactúe con la arquitectura interna ni seleccione métodos de \textit{parsing}. Esta abstracción de la complejidad técnica simplifica la experiencia de uso y facilita la adopción del producto en entornos operativos. Complementariamente, el cliente cuenta con acceso a un \textit{dashboard} donde puede monitorear métricas de desempeño, como el porcentaje de error de los resultados obtenidos, y además dispone de funcionalidades para configurar o crear sus propios extractores basados en \textit{LLM} según sus necesidades específicas.

\section{Arquitectura de la solución propuesta}

El flujo de producción representa el proceso operativo que sigue el sistema ante una solicitud real del usuario final. Como se observa en la figura~\ref{fig:API-interactor}, la interacción comienza en el entorno web, desde donde el usuario carga un documento en formato PDF, JPG o PNG. Una vez recibido, el archivo es almacenado en S3 y posteriormente recuperado por el \textit{backend} para su procesamiento. El sistema ejecuta entonces el proceso de extracción mediante un extractor basado en LLM, el cual opera a partir de una configuración compuesta por el \textit{prompt}, el modelo y el esquema de salida esperado. El modelo interpreta el contenido del documento y genera una salida estructurada en formato JSON con los campos extraídos. Finalmente, el entorno web presenta los resultados en campos editables que permiten al usuario revisar y corregir la información antes de confirmar el registro. Los detalles de implementación, como el preprocesamiento de imágenes y los mecanismos de validación y reintento, se describen en el capítulo~\ref{chap:mvp}.


\begin{figure}[H]
\centering
\resizebox{0.45\textwidth}{!}{
\begin{tikzpicture}[
    node distance=1cm,
    every node/.style={font=\small, align=center},
    entrada/.style={
        rectangle, rounded corners=6pt, draw=gray!60, thick,
        minimum width=6.5cm, minimum height=1.4cm, fill=gray!8
    },
    proceso/.style={
        rectangle, rounded corners=6pt, draw=orange!50, thick,
        minimum width=6.5cm, minimum height=1.4cm, fill=orange!8
    },
    resultado/.style={
        rectangle, rounded corners=6pt, draw=teal!50, thick,
        minimum width=6.5cm, minimum height=1.4cm, fill=teal!8
    },
    config/.style={
        rectangle, rounded corners=6pt, draw=cyan!40, thick,
        minimum width=3.8cm, minimum height=1.4cm, fill=cyan!10
    },
    etiqueta/.style={
        font=\scriptsize\itshape,
        text=gray!70
    },
    flecha/.style={
        -{Latex[length=3mm]},
        thick
    }
]

% Nodos verticales
\node[entrada] (usuario) {
    Usuario\\[1pt]
    {\scriptsize PDF / JPG / PNG}
};
\node[entrada, below=of usuario] (web) {
    Entorno web
};
\node[proceso, below=of web] (backend) {
    \textit{Backend} (FastAPI)\\[1pt]
    {\scriptsize Almacenamiento S3 $\cdot$ Orquestación}
};
\node[proceso, below=of backend] (extractor) {
    Extractor basado en \textit{LLM}\\[1pt]
    {\scriptsize visión + salida estructurada (JSON)}
};
\node[resultado, below=of extractor] (registro) {
    Registro validado
};

% Config lateral
\node[config, right=1.5cm of extractor] (confignode) {
    Extractor \textit{config}\\[1pt]
    {\scriptsize \textit{prompt} + modelo + \textit{schema}}
};

% Flechas verticales
\draw[flecha] (usuario) -- node[right, etiqueta] {} (web);
\draw[flecha] (web) -- node[right, etiqueta] {\textit{presigned upload} + API \textit{call}} (backend);
\draw[flecha] (backend) -- node[right, etiqueta] {documento + \textit{config}} (extractor);
\draw[flecha] (extractor) -- node[right, etiqueta] {revisión y confirmación} (registro);

% Flecha config
\draw[flecha, dashed, gray!60] (confignode) -- (extractor);

\end{tikzpicture}
}
\caption{Interacción del usuario con la plataforma}
\label{fig:API-interactor}
\end{figure}

\section{Requisitos funcionales y no funcionales mínimos}

En esta sección se presenta el conjunto mínimo de requisitos que orienta el funcionamiento del sistema de extracción y evaluación. Se distinguen, por un lado, las capacidades específicas que el sistema debe ofrecer y, por otro lado, las condiciones de calidad que deben cumplirse de forma constante. La Tabla~\ref{tab:requisitos_minimos} detalla estos requisitos. Su definición permite asegurar un comportamiento estable del proceso a medida que se incorporan nuevos documentos y extractores.

\begin{table}[htbp]
\centering
\resizebox{\textwidth}{!}{%
\begin{tabular}{p{1.2cm} p{10.0cm} p{1.2cm} p{10.0cm}}
\toprule
\multicolumn{2}{c}{\textbf{Requisitos funcionales}} & \multicolumn{2}{c}{\textbf{Requisitos no funcionales}} \\
\midrule
\textbf{Código} & \textbf{Descripción} & \textbf{Código} & \textbf{Descripción} \\
\midrule
RF1 & Ingesta de documentos: aceptar archivos CSV con referencias a PDFs, descargarlos y validarlos. &
RNF1 & Precisión: alcanzar un mínimo de 85\% de precisión en los tres campos sobre los documentos de prueba. \\
RF2 & Validación de datos: verificar que las CLABEs tengan 18 dígitos, que los RFC cumplan el formato oficial mexicano y que las CURP tengan 18 caracteres con estructura válida. &
RNF2 & Latencia: responder en menos de 10 segundos para documentos de hasta 5 páginas. \\
RF3 & Extracción de campos: extraer correctamente los campos principales definidos para el análisis. &
RNF3 & Disponibilidad: mantener una disponibilidad mínima de 99\% durante el horario laboral. \\
RF4 & Comparación de métodos: ejecutar múltiples extractores sobre los mismos documentos y generar reportes comparativos de precisión. &
RNF4 & Escalabilidad: soportar al menos 100 documentos por hora. \\
RF5 & API REST: exponer \textit{endpoints} HTTP que acepten archivos en formato PDF, JPG o PNG y retornen JSON con los campos extraídos. &
RNF5 & Reproducibilidad: todas las evaluaciones deben ser reproducibles. \\
RF6 & \textit{Logging}: registrar cada operación en \textit{logs} estructurados con nivel de severidad. &
RNF6 & Mantenibilidad: el código debe seguir principios SOLID. \\
\bottomrule
\end{tabular}%
}
\caption{Requisitos funcionales y no funcionales mínimos del sistema}
\label{tab:requisitos_minimos}
\end{table}

    